\makeatletter
\def\input@path{{./}{../}{../../}{../../../}{../../../../}{preamble/}{../preamble/}{../../preamble/}{../../../preamble/}{../../../../preamble/}{book2026-2027/preamble/}{book2026-2027/}}
\makeatother

\input{preamble_G10_T1_C1C2C3}
\input{preamble_exam_layout}

\title{Decision-Making Under Uncertainty: Exam Solutions}
\author{}
\date{}

\begin{document}
	\maketitle
	\vspace{-2cm}
	\setcounter{tocdepth}{1}
	\setcounter{secnumdepth}{2}
	\phantomsection
	\label{toc}
\label{sec:solutions-extraction}
	\tableofcontents
	
	\vspace{-5mm}
	\section*{Learning objective}\vspace{-5mm}
	Structure decision-making problems under uncertainty, calculate and organise payoffs, and apply the Maximax and Maximin criteria to compare alternatives and justify decisions.
	
	\section*{Criteria assessment}\vspace{-5mm}
	This task sheet assesses criteria C1 through C4. Each problem assesses the criteria indicated in its title. The number of points assigned to each task is shown at the end of the item. Partial credit may be awarded when the correct procedure is shown but arithmetic errors are made. Answers that do not show the required procedure will not receive credit. A minimum of 4 points is required to pass each criterion.
	
	\section*{Evaluated criteria}
	\begin{enumerate}[label=\textbf{C\arabic*}]
		\item Characterises decision-making problems under uncertainty by identifying decision alternatives, states of nature, consequences, and the economic or financial context of the decision.
		\item Organises decision problems using payoff tables, placing alternatives, states of nature, and payoffs in a clear structure for comparison.
		\item Calculates profits by combining revenues with fixed and variable costs across different states of nature, determining the corresponding payoffs for each alternative, and organising the results in payoff tables for comparison.
		\item Applies the Maximin and Maximax criteria to make optimistic or conservative decisions without probabilities, selecting the alternative with the highest possible payoff or the best worst-case, interpreting risks or limitations.
	\end{enumerate}
	
	\sectionScored
	{Meal Delivery Service}
	{C1-\sml{1} \hspace{2mm} C2-\sml{2}}
	{sec:meal-delivery-service-c1-c2}
	{
		A meal delivery company is planning its service offering for an upcoming operating period. A standard delivery service is expected to generate profits of twenty thousand US dollars under high demand and twelve thousand US dollars under low demand, while a premium delivery service is expected to generate thirty-four thousand US dollars and seven thousand US dollars, respectively. The company must select its delivery service before the actual level of customer demand is known.
		
		\textbf{Questions/Tasks.}
		\begin{enumerate}
			\CriterionTask
			{subsec:meal-delivery-service-c1-c2-c1-characterization}{1}
			{Characterize the key elements of the decision-making situation. \score{1}}
			
			\CriterionTask
			{subsec:meal-delivery-service-c1-c2-c2-payoff-table}{2}
			{Construct the payoff table. \score{2}}
		\end{enumerate}
		
		\phantomsection
		\label{subsec:meal-delivery-service-c1-c2-c1-characterization}
		\subsection*{C1 - Characterization \score{1}}
		
		\DecisionElementsTable{Standard delivery service; Premium delivery service}{High demand; Low demand}{Actual customer demand after the delivery service decision}{Operating the selected activity under the realised conditions leaves funds available or a shortfall to cover. Monetary measure: Profit in thousand US dollars for each alternative--state pair}
		
		\phantomsection
		\label{subsec:meal-delivery-service-c1-c2-c2-payoff-table}
		\subsection*{C2 - Payoff table \score{2}}
		
		Each supplied monetary payoff is placed at the intersection of its alternative and state:
		\[
		\begin{array}{ll}
			P(A_1,S_1)=20, & P(A_1,S_2)=12,\\
			P(A_2,S_1)=34, & P(A_2,S_2)=7.
		\end{array}
		\]
		Thus, in thousand US dollars, the completed payoff table is
		\begin{center}
			\begin{tabular}{lcc}
				\toprule
				Alternative & High \(S_1\) & Low \(S_2\)\\
				\midrule
				Standard delivery service \(A_1\) & 20 & 12\\
				Premium delivery service \(A_2\) & 34 & 7\\
				\bottomrule
			\end{tabular}
		\end{center}
		
		\vspace{-5mm}
		\secInicio[sec:meal-delivery-service-c1-c2][subsec:meal-delivery-service-c1-c2-c2-payoff-table]
		
	}
	
	\sectionScored
	{Small Coffee Roastery Production}
	{C1-\sml{1} \hspace{2mm} C2-\sml{3}}
	{sec:small-coffee-roastery-production-c1-c2}
	{
		A small coffee roastery is planning its production capacity for an upcoming sales period. A standard roasting machine can generate 4 thousand BRL per production batch, while an industrial roasting machine can generate 6 thousand BRL per batch. Based on expected demand, the standard roasting machine would complete 18 batches under high demand or 12 under low demand, whereas the industrial roasting machine would complete 15 or 10 batches, respectively. In addition, the roastery would receive a fixed sales bonus of 7 thousand BRL if demand is high and incur a fixed loss of 4 thousand BRL if demand is low.
		
		\textbf{Questions/Tasks.}
		\begin{enumerate}
			\CriterionTask
			{subsec:small-coffee-roastery-production-c1-c2-c1-characterization}{1}
			{Characterize the key elements of the decision-making situation. \score{1}}
			
			\CriterionTask
			{subsec:small-coffee-roastery-production-c1-c2-c2-payoff-table}{2}
			{Construct the payoff table. \score{3}}
		\end{enumerate}
		
		\phantomsection
		\label{subsec:small-coffee-roastery-production-c1-c2-c1-characterization}
		\subsection*{C1 - Characterization \score{1}}
		
		\DecisionElementsTable
{Standard roasting machine; Industrial roasting machine}
{High demand; Low demand}
{Actual customer demand after the roasting machine decision}
{Operating the selected activity under the realised conditions leaves funds available or a shortfall to cover. The monetary measure is profit in thousand BRL for each alternative--state pair.}

		
		\begin{CalculationLandscape}{11in}{8.5in}
\phantomsection
		\label{subsec:small-coffee-roastery-production-c1-c2-c2-payoff-table}
		\subsection*{C2 - Payoff table \score{3}}
		
		
\SymbolsTable{
$A_i$ & $A_1$: Standard roasting machine; $A_2$: Industrial roasting machine & Alternative labels, in row order. \\
$S_j$ & $S_1$: High demand; $S_2$: Low demand & State labels, in column order; revealed after the choice. \\
$q(A_i,S_j)$ & Standard--high: 18 batches; Standard--low: 12 batches; Industrial--high: 15 batches; Industrial--low: 10 batches & Expected completed production batches for each roasting machine under each demand condition. \\
$v^A(A_i)$ & Standard roasting machine: 4 thousand BRL/batch; Industrial roasting machine: 6 thousand BRL/batch & Earnings per completed production batch from the chosen roasting machine. \\
$f^S(S_j)$ & High demand: 7 thousand BRL; Low demand: $-4$ thousand BRL & Fixed sales bonus or loss associated with the demand condition. \\
}

		
		\textbf{Payoff equation.}
		\[
		P_{ay}(A_i,S_j)
		=
		q(A_i,S_j)v^A(A_i)
		+
		f^S(S_j).
		\]
		
		\CalculationPair{Additive payoff}{
\begin{tabular}{@{}lcc@{}}
\toprule Alternative & $S_1$ & $S_2$ \\ \midrule
$A_1$ & \SymbolicCell{q(A_1,S_1)\times v^A(A_1) \\ {}+f^S(S_1)}{P_{ay}(A_1,S_1)} & \SymbolicCell{q(A_1,S_2)\times v^A(A_1) \\ {}+f^S(S_2)}{P_{ay}(A_1,S_2)} \\[5pt]
$A_2$ & \SymbolicCell{q(A_2,S_1)\times v^A(A_2) \\ {}+f^S(S_1)}{P_{ay}(A_2,S_1)} & \SymbolicCell{q(A_2,S_2)\times v^A(A_2) \\ {}+f^S(S_2)}{P_{ay}(A_2,S_2)} \\[5pt]
\bottomrule
\end{tabular}
}{
\begin{tabular}{@{}lcc@{}}
\toprule Alternative & $S_1$ & $S_2$ \\ \midrule
$A_1$ & \CalcResult{18\times 4 \\ {}+7}{79} & \CalcResult{12\times 4 \\ {}-4}{44} \\[5pt]
$A_2$ & \CalcResult{15\times 6 \\ {}+7}{97} & \CalcResult{10\times 6 \\ {}-4}{56} \\[5pt]
\bottomrule
\end{tabular}
}
Negative rates reduce the per-unit contribution; negative fixed amounts represent losses or commitments deducted once. A negative final payoff is a shortfall.

		
		\begingroup\renewcommand{\CalculationStateCount}{2}
\FinalPayoffTable{profit in thousand BRL}{High demand ($S_1$) & Low demand ($S_2$)}{
Standard roasting machine ($A_1$) & 79 & 44 \\
Industrial roasting machine ($A_2$) & 97 & 56 \\
}
\endgroup
\end{CalculationLandscape}

		
		\vspace{-5mm}
		\secInicio[sec:small-coffee-roastery-production-c1-c2][subsec:small-coffee-roastery-production-c1-c2-c2-payoff-table]
	}
	
	
	\sectionScored
	{Online Fitness Membership}
	{C1-\sml{2} \hspace{2mm} C2-\sml{5}}
	{sec:online-fitness-membership-c1-c2}
	{
		An online fitness company must choose between a standard membership and a flexible membership before customer demand is known. Demand may be high or low. The standard membership generates 5 thousand BRL per active member and also provides a fixed operating contribution of 5 thousand BRL. The flexible membership generates 7 thousand BRL per active member and provides a fixed operating contribution of 8 thousand BRL. High demand adds another 3 thousand BRL per active member to either membership, while low demand adds 1 thousand BRL per active member. The standard membership is expected to have 15 active members under high demand and 9 under low demand; the flexible membership is expected to have 12 active members under high demand and 7 under low demand.
		
		\textbf{Questions/Tasks.}
		\begin{enumerate}
			\CriterionTask
			{subsec:online-fitness-membership-c1-c2-c1-characterization}{1}
			{Characterize the key elements of the decision-making situation. \score{2}}
			
			\CriterionTask
			{subsec:online-fitness-membership-c1-c2-c2-payoff-table}{2}
			{Construct the payoff table. \score{5}}
		\end{enumerate}
		
		\phantomsection
		\label{subsec:online-fitness-membership-c1-c2-c1-characterization}
		\subsection*{C1 - Characterization \score{2}}
		
		\DecisionElementsTable
{Standard membership; Flexible membership}
{High customer demand; Low customer demand}
{Actual customer demand after the membership is selected}
{Operating the selected activity under the realised conditions leaves funds available or a shortfall to cover. The monetary measure is profit in thousand BRL for each alternative--state pair.}

		
		\begin{CalculationLandscape}{11in}{8.5in}
\phantomsection
		\label{subsec:online-fitness-membership-c1-c2-c2-payoff-table}
		\subsection*{C2 - Payoff table \score{5}}
		
		
\SymbolsTable{
$A_i$ & $A_1$: Standard membership; $A_2$: Flexible membership & Alternative labels, in row order. \\
$S_j$ & $S_1$: High customer demand; $S_2$: Low customer demand & State labels, in column order; revealed after the choice. \\
$q(A_i,S_j)$ & Standard--high: 15 members; Standard--low: 9 members; Flexible--high: 12 members; Flexible--low: 7 members & Expected active members for each membership under each level of customer demand. \\
$v^A(A_i)$ & Standard membership: 5 thousand BRL/member; Flexible membership: 7 thousand BRL/member & Contribution per active member from the chosen membership. \\
$v^S(S_j)$ & High demand: 3 thousand BRL/member; Low demand: 1 thousand BRL/member & Additional contribution per active member from customer demand. \\
$f^A(A_i)$ & Standard membership: 5 thousand BRL; Flexible membership: 8 thousand BRL & Fixed operating contribution from the chosen membership. \\
}

		
		\textbf{Payoff equation.}
		\[
		P_{ay}(A_i,S_j)
		=
		q(A_i,S_j)
		\left[
		v^A(A_i)+v^S(S_j)
		\right]
		+
		f^A(A_i).
		\]
		
		\CalculationPair{Additive payoff}{
\begin{tabular}{@{}lcc@{}}
\toprule Alternative & $S_1$ & $S_2$ \\ \midrule
$A_1$ & \SymbolicCell{q(A_1,S_1)\times \\ [v^A(A_1)+v^S(S_1)] \\ {}+f^A(A_1)}{P_{ay}(A_1,S_1)} & \SymbolicCell{q(A_1,S_2)\times \\ [v^A(A_1)+v^S(S_2)] \\ {}+f^A(A_1)}{P_{ay}(A_1,S_2)} \\[5pt]
$A_2$ & \SymbolicCell{q(A_2,S_1)\times \\ [v^A(A_2)+v^S(S_1)] \\ {}+f^A(A_2)}{P_{ay}(A_2,S_1)} & \SymbolicCell{q(A_2,S_2)\times \\ [v^A(A_2)+v^S(S_2)] \\ {}+f^A(A_2)}{P_{ay}(A_2,S_2)} \\[5pt]
\bottomrule
\end{tabular}
}{
\begin{tabular}{@{}lcc@{}}
\toprule Alternative & $S_1$ & $S_2$ \\ \midrule
$A_1$ & \CalcResult{15\times \\ [5+3] \\ {}+5}{125} & \CalcResult{9\times \\ [5+1] \\ {}+5}{59} \\[5pt]
$A_2$ & \CalcResult{12\times \\ [7+3] \\ {}+8}{128} & \CalcResult{7\times \\ [7+1] \\ {}+8}{64} \\[5pt]
\bottomrule
\end{tabular}
}

		
		\begingroup\renewcommand{\CalculationStateCount}{2}
\FinalPayoffTable{profit in thousand BRL}{High customer demand ($S_1$) & Low customer demand ($S_2$)}{
Standard membership ($A_1$) & 125 & 59 \\
Flexible membership ($A_2$) & 128 & 64 \\
}
\endgroup
\end{CalculationLandscape}

		
		\vspace{-5mm}
		\secInicio[sec:online-fitness-membership-c1-c2][subsec:online-fitness-membership-c1-c2-c2-payoff-table]
		
	}
	
	\sectionScored
	{Neighborhood Bakery Expansion}
	{C1-\sml{1} \hspace{2mm} C3-\sml{2}}
	{sec:neighborhood-bakery-expansion-c1-c3}
	{
		A neighborhood bakery must decide whether to introduce a breakfast menu or maintain its current product selection before customer traffic is known. Traffic may be strong or weak. All values are in million BRL. Introducing the breakfast menu produces revenue of 68 under strong traffic and 38 under weak traffic, while maintaining the current selection produces revenue of 59 under strong traffic and 49 under weak traffic. The corresponding costs are 23 and 21 for the breakfast menu under strong and weak traffic, respectively, and 24 and 22 for the current selection under strong and weak traffic, respectively.
			\textbf{Questions/Tasks.}
		\begin{enumerate}
			\CriterionTask
			{subsec:neighborhood-bakery-expansion-c1-c3-c1-characterization}{1}
			{Characterize the key elements of the decision-making situation. \score{1}}
			
			\CriterionTask
			{subsec:neighborhood-bakery-expansion-c1-c3-c3-profit-payoff-table}{3}
			{Use revenue and costs to construct the payoff table. \score{2}}
		\end{enumerate}
		
		\phantomsection
		\label{subsec:neighborhood-bakery-expansion-c1-c3-c1-characterization}
		\subsection*{C1 - Characterization \score{1}}
		
		\DecisionElementsTable
{Introduce breakfast menu; Maintain current selection}
{Strong traffic; Weak traffic}
{Actual customer traffic after the product selection decision}
{Operating the selected activity under the realised conditions leaves funds available or a shortfall to cover. The monetary measure is profit in million BRL for each alternative--state pair.}

		
		\begin{CalculationLandscape}{11in}{8.5in}
\phantomsection
		\label{subsec:neighborhood-bakery-expansion-c1-c3-c3-profit-payoff-table}
		\subsection*{C3 - Calculating profit and building the payoff table \score{2}}
		
		
\SymbolsTable{
$A_i$ & $A_1$: Introduce breakfast menu; $A_2$: Maintain current selection & Alternative labels, in row order. \\
$S_j$ & $S_1$: Strong traffic; $S_2$: Weak traffic & State labels, in column order; revealed after the choice. \\
$R(A_i,S_j)$ & Breakfast menu--strong: 68 million BRL; Breakfast menu--weak: 38 million BRL; Current selection--strong: 59 million BRL; Current selection--weak: 49 million BRL & Total revenue generated by each product selection under each traffic condition. \\
$C(A_i,S_j)$ & Breakfast menu--strong: 23 million BRL; Breakfast menu--weak: 21 million BRL; Current selection--strong: 24 million BRL; Current selection--weak: 22 million BRL & Total cost associated with each product selection under each traffic condition. \\
}

		
		\textbf{Profit equation.}
		\[
		P(A_i,S_j)
		=
		R(A_i,S_j)-C(A_i,S_j).
		\]
		
		\CalculationPair{Profit}{
\begin{tabular}{@{}lcc@{}}
\toprule Alternative & $S_1$ & $S_2$ \\ \midrule
$A_1$ & \SymbolicCell{R(A_1,S_1)-C(A_1,S_1)}{P(A_1,S_1)} & \SymbolicCell{R(A_1,S_2)-C(A_1,S_2)}{P(A_1,S_2)} \\[5pt]
$A_2$ & \SymbolicCell{R(A_2,S_1)-C(A_2,S_1)}{P(A_2,S_1)} & \SymbolicCell{R(A_2,S_2)-C(A_2,S_2)}{P(A_2,S_2)} \\[5pt]
\bottomrule
\end{tabular}
}{
\begin{tabular}{@{}lcc@{}}
\toprule Alternative & $S_1$ & $S_2$ \\ \midrule
$A_1$ & \CalcResult{68-23}{45} & \CalcResult{38-21}{17} \\[5pt]
$A_2$ & \CalcResult{59-24}{35} & \CalcResult{49-22}{27} \\[5pt]
\bottomrule
\end{tabular}
}

		
		\begingroup\renewcommand{\CalculationStateCount}{2}
\FinalPayoffTable{profit in million BRL}{Strong traffic ($S_1$) & Weak traffic ($S_2$)}{
Introduce breakfast menu ($A_1$) & 45 & 17 \\
Maintain current selection ($A_2$) & 35 & 27 \\
}
\endgroup
\end{CalculationLandscape}

		
		\vspace{-5mm}
		\secInicio[sec:neighborhood-bakery-expansion-c1-c3][subsec:neighborhood-bakery-expansion-c1-c3-c3-profit-payoff-table]
		
	}
	
	\sectionScored
	{Bakery Product Production}
	{C1-\sml{1} \hspace{2mm} C3-\sml{3}}
	{sec:bakery-products-c1-c3}
	{
		A bakery is planning its production for an upcoming holiday season. Revenue per production batch is 20 thousand EUR for cakes and
		16 thousand EUR for pastries. Under high customer demand, the
		bakery expects to produce 9 batches of cakes and 6 batches of pastries, while
		under low customer demand it expects to produce 5 batches of
		cakes and 4 batches of pastries. The cost per production batch is 8 thousand
		EUR under high customer demand and 11 thousand EUR under low
		customer demand.
		
			\textbf{Questions/Tasks.}
		\begin{enumerate}
			\CriterionTask
			{subsec:bakery-products-c1-c3-c1-characterization}{1}
			{Characterize the key elements of the decision-making situation. \score{1}}
			
			\CriterionTask
			{subsec:bakery-products-c1-c3-c3-profit-table}{3}
			{Use revenue and costs to construct the payoff table. \score{3}}
		\end{enumerate}
		
		\phantomsection
		\label{subsec:bakery-products-c1-c3-c1-characterization}
		\subsection*{C1 - Characterization \score{1}}
		
		\DecisionElementsTable
{Produce cakes; Produce pastries}
{High customer demand; Low customer demand}
{Customer demand revealed after the production decision}
{Operating the selected activity under the realised conditions leaves funds available or a shortfall to cover. The monetary measure is profit in thousand EUR for each product--demand-state pair.}

		
		\begin{CalculationLandscape}{14in}{11in}
\phantomsection
		\label{subsec:bakery-products-c1-c3-c3-profit-table}
		\subsection*{C3 - Profit and payoff table \score{3}}
		
		\secInicio[sec:solutions-extraction][sec:bakery-products-c1-c3]
		

\SymbolsTable{
$A_i$ & $A_1$: Produce cakes; $A_2$: Produce pastries & Alternative labels, in row order. \\
$S_j$ & $S_1$: High customer demand; $S_2$: Low customer demand & State labels, in column order; revealed after the choice. \\
$q(A_i,S_j)$ & Cakes--high: 9; Cakes--low: 5; Pastries--high: 6; Pastries--low: 4 & Production batches for each product under each demand condition. \\
$r(A_i)$ & Cakes: 20 thousand EUR/batch; Pastries: 16 thousand EUR/batch & Revenue per production batch. \\
$c(S_j)$ & High customer demand: 8 thousand EUR/batch; Low customer demand: 11 thousand EUR/batch & Cost per production batch. \\
}

		
		\textbf{Profit equation.}
		\[
		R=q(A_i,S_j)r(A_i),
		\qquad
		C=q(A_i,S_j)c(S_j),
		\qquad
		P(A_i,S_j)=R-C.
		\]
		
		\CalculationPair{Revenue}{
\begin{tabular}{@{}lcc@{}}
\toprule Alternative & $S_1$ & $S_2$ \\ \midrule
$A_1$ & \SymbolicCell{q(A_1,S_1)\times r(A_1)}{R(A_1,S_1)} & \SymbolicCell{q(A_1,S_2)\times r(A_1)}{R(A_1,S_2)} \\[5pt]
$A_2$ & \SymbolicCell{q(A_2,S_1)\times r(A_2)}{R(A_2,S_1)} & \SymbolicCell{q(A_2,S_2)\times r(A_2)}{R(A_2,S_2)} \\[5pt]
\bottomrule
\end{tabular}
}{
\begin{tabular}{@{}lcc@{}}
\toprule Alternative & $S_1$ & $S_2$ \\ \midrule
$A_1$ & \CalcResult{9\times 20}{180} & \CalcResult{5\times 20}{100} \\[5pt]
$A_2$ & \CalcResult{6\times 16}{96} & \CalcResult{4\times 16}{64} \\[5pt]
\bottomrule
\end{tabular}
}
\CalculationPair{Cost}{
\begin{tabular}{@{}lcc@{}}
\toprule Alternative & $S_1$ & $S_2$ \\ \midrule
$A_1$ & \SymbolicCell{q(A_1,S_1)\times c(S_1)}{C(A_1,S_1)} & \SymbolicCell{q(A_1,S_2)\times c(S_2)}{C(A_1,S_2)} \\[5pt]
$A_2$ & \SymbolicCell{q(A_2,S_1)\times c(S_1)}{C(A_2,S_1)} & \SymbolicCell{q(A_2,S_2)\times c(S_2)}{C(A_2,S_2)} \\[5pt]
\bottomrule
\end{tabular}
}{
\begin{tabular}{@{}lcc@{}}
\toprule Alternative & $S_1$ & $S_2$ \\ \midrule
$A_1$ & \CalcResult{9\times 8}{72} & \CalcResult{5\times 11}{55} \\[5pt]
$A_2$ & \CalcResult{6\times 8}{48} & \CalcResult{4\times 11}{44} \\[5pt]
\bottomrule
\end{tabular}
}
\CalculationPair{Profit}{
\begin{tabular}{@{}lcc@{}}
\toprule Alternative & $S_1$ & $S_2$ \\ \midrule
$A_1$ & \SymbolicCell{R(A_1,S_1)-C(A_1,S_1)}{P(A_1,S_1)} & \SymbolicCell{R(A_1,S_2)-C(A_1,S_2)}{P(A_1,S_2)} \\[5pt]
$A_2$ & \SymbolicCell{R(A_2,S_1)-C(A_2,S_1)}{P(A_2,S_1)} & \SymbolicCell{R(A_2,S_2)-C(A_2,S_2)}{P(A_2,S_2)} \\[5pt]
\bottomrule
\end{tabular}
}{
\begin{tabular}{@{}lcc@{}}
\toprule Alternative & $S_1$ & $S_2$ \\ \midrule
$A_1$ & \CalcResult{180-72}{108} & \CalcResult{100-55}{45} \\[5pt]
$A_2$ & \CalcResult{96-48}{48} & \CalcResult{64-44}{20} \\[5pt]
\bottomrule
\end{tabular}
}

		
		\begingroup\renewcommand{\CalculationStateCount}{2}
\FinalPayoffTable{profit in thousand EUR}{High customer demand ($S_1$) & Low customer demand ($S_2$)}{
Produce cakes ($A_1$) & 108 & 45 \\
Produce pastries ($A_2$) & 48 & 20 \\
}
\endgroup
\end{CalculationLandscape}

		
		
	}
		
	\sectionScored
	{Meal Delivery Service Packages}
	{C1-\sml{2} \hspace{2mm} C3-\sml{5}}
	{sec:meal-delivery-c1-c3}
	{
		A meal delivery service must choose between a standard meal package and a
		premium meal package before weekly demand is known. Demand may be
		high or low. Quantities, in hundreds of meal packages, are 10 and 7
		for the standard package and 8 and 5 for the premium package under high and
		low demand, respectively. The package-based revenue rates are 7 and 12
		thousand USD per hundred meal packages for the standard and premium
		packages. High demand adds a state-based revenue rate of 3 thousand USD
		per hundred meal packages; low demand adds 1. The package-based cost rates
		are 3 and 5 thousand USD per hundred meal packages, respectively. The
		fixed operating costs are 14 thousand USD under high demand and
		18 thousand USD under low demand.
		
			\textbf{Questions/Tasks.}
		\begin{enumerate}
			\CriterionTask
			{subsec:meal-delivery-c1-c3-c1-characterization}{1}
			{Characterize the key elements of the decision-making situation. \score{2}}
			
			\CriterionTask
			{subsec:meal-delivery-c1-c3-c3-profit-table}{3}
			{Use revenue and costs to construct the payoff table. \score{5}}
		\end{enumerate}
		
		\phantomsection
		\label{subsec:meal-delivery-c1-c3-c1-characterization}
		\subsection*{C1 - Characterization \score{2}}
		
		\DecisionElementsTable
{Offer a standard meal package; Offer a premium meal package}
{High weekly demand; Low weekly demand}
{Weekly market conditions revealed after the meal-package decision}
{Operating the selected activity under the realised conditions leaves funds available or a shortfall to cover. The monetary measure is profit in thousand USD for each package--market-state pair.}

		
		\begin{CalculationLandscape}{14in}{11in}
\phantomsection
		\label{subsec:meal-delivery-c1-c3-c3-profit-table}
		\subsection*{C3 - Profit and payoff table \score{5}}
		
		\secInicio[sec:solutions-extraction][sec:meal-delivery-c1-c3]
		

\SymbolsTable{
$A_i$ & $A_1$: Offer a standard meal package; $A_2$: Offer a premium meal package & Alternative labels, in row order. \\
$S_j$ & $S_1$: High weekly demand; $S_2$: Low weekly demand & State labels, in column order; revealed after the choice. \\
$q(A_i,S_j)$ & Standard--high: 10; Standard--low: 7; Premium--high: 8; Premium--low: 5 hundreds of meal packages & Meal-package lots for each package under each demand level. \\
$r^A(A_i)$ & Standard package: 7 thousand USD/100 meal packages; Premium package: 12 thousand USD/100 meal packages & Package-based revenue rate per hundred meal packages. \\
$r^S(S_j)$ & High demand: 3 thousand USD/100 meal packages; Low demand: 1 thousand USD/100 meal packages & Additional revenue rate per hundred meal packages from weekly demand. \\
$c^A(A_i)$ & Standard package: 3 thousand USD/100 meal packages; Premium package: 5 thousand USD/100 meal packages & Package-based cost rate per hundred meal packages. \\
$f_c(S_j)$ & High demand: 14 thousand USD; Low demand: 18 thousand USD & Fixed operating cost under each demand condition. \\
}

		
		\textbf{Profit equation.}
		\[
		R=q(A_i,S_j)\left[r^A(A_i)+r^S(S_j)\right],
		\qquad
		C=q(A_i,S_j)c^A(A_i)+f_c(S_j),
		\qquad
		P(A_i,S_j)=R-C.
		\]
		
		\CalculationPair{Revenue}{
\begin{tabular}{@{}lcc@{}}
\toprule Alternative & $S_1$ & $S_2$ \\ \midrule
$A_1$ & \SymbolicCell{q(A_1,S_1)\times \\ [r^A(A_1)+r^S(S_1)]}{R(A_1,S_1)} & \SymbolicCell{q(A_1,S_2)\times \\ [r^A(A_1)+r^S(S_2)]}{R(A_1,S_2)} \\[5pt]
$A_2$ & \SymbolicCell{q(A_2,S_1)\times \\ [r^A(A_2)+r^S(S_1)]}{R(A_2,S_1)} & \SymbolicCell{q(A_2,S_2)\times \\ [r^A(A_2)+r^S(S_2)]}{R(A_2,S_2)} \\[5pt]
\bottomrule
\end{tabular}
}{
\begin{tabular}{@{}lcc@{}}
\toprule Alternative & $S_1$ & $S_2$ \\ \midrule
$A_1$ & \CalcResult{10\times \\ [7+3]}{100} & \CalcResult{7\times \\ [7+1]}{56} \\[5pt]
$A_2$ & \CalcResult{8\times \\ [12+3]}{120} & \CalcResult{5\times \\ [12+1]}{65} \\[5pt]
\bottomrule
\end{tabular}
}
\CalculationPair{Cost}{
\begin{tabular}{@{}lcc@{}}
\toprule Alternative & $S_1$ & $S_2$ \\ \midrule
$A_1$ & \SymbolicCell{q(A_1,S_1)\times c^A(A_1) \\ {}+f_c(S_1)}{C(A_1,S_1)} & \SymbolicCell{q(A_1,S_2)\times c^A(A_1) \\ {}+f_c(S_2)}{C(A_1,S_2)} \\[5pt]
$A_2$ & \SymbolicCell{q(A_2,S_1)\times c^A(A_2) \\ {}+f_c(S_1)}{C(A_2,S_1)} & \SymbolicCell{q(A_2,S_2)\times c^A(A_2) \\ {}+f_c(S_2)}{C(A_2,S_2)} \\[5pt]
\bottomrule
\end{tabular}
}{
\begin{tabular}{@{}lcc@{}}
\toprule Alternative & $S_1$ & $S_2$ \\ \midrule
$A_1$ & \CalcResult{10\times 3 \\ {}+14}{44} & \CalcResult{7\times 3 \\ {}+18}{39} \\[5pt]
$A_2$ & \CalcResult{8\times 5 \\ {}+14}{54} & \CalcResult{5\times 5 \\ {}+18}{43} \\[5pt]
\bottomrule
\end{tabular}
}
\CalculationPair{Profit}{
\begin{tabular}{@{}lcc@{}}
\toprule Alternative & $S_1$ & $S_2$ \\ \midrule
$A_1$ & \SymbolicCell{R(A_1,S_1)-C(A_1,S_1)}{P(A_1,S_1)} & \SymbolicCell{R(A_1,S_2)-C(A_1,S_2)}{P(A_1,S_2)} \\[5pt]
$A_2$ & \SymbolicCell{R(A_2,S_1)-C(A_2,S_1)}{P(A_2,S_1)} & \SymbolicCell{R(A_2,S_2)-C(A_2,S_2)}{P(A_2,S_2)} \\[5pt]
\bottomrule
\end{tabular}
}{
\begin{tabular}{@{}lcc@{}}
\toprule Alternative & $S_1$ & $S_2$ \\ \midrule
$A_1$ & \CalcResult{100-44}{56} & \CalcResult{56-39}{17} \\[5pt]
$A_2$ & \CalcResult{120-54}{66} & \CalcResult{65-43}{22} \\[5pt]
\bottomrule
\end{tabular}
}

		
		\begingroup\renewcommand{\CalculationStateCount}{2}
\FinalPayoffTable{profit in thousand USD}{High weekly demand ($S_1$) & Low weekly demand ($S_2$)}{
Standard meal package ($A_1$) & 56 & 17 \\
Premium meal package ($A_2$) & 66 & 22 \\
}
\endgroup
\end{CalculationLandscape}

		
	}
		
	\sectionScored
	{Solar Energy Investment Strategy}
	{C4-\sml{2}}
	{sec:solar-energy-investment-c4}
	{
		A renewable energy company must choose a strategy for investing in a new
		market. Energy demand could be high, medium, or low. All payoffs are in
		millions of dollars.
		
			\begin{center}
			\renewcommand{\arraystretch}{1.3}
			\small
			\begin{tabular}{lccc}
				\toprule
				& \multicolumn{3}{c}{Energy Demand} \\
				\cmidrule(lr){2-4}
				Alternative & High & Medium & Low \\
				\midrule
				Utility-scale solar farm & 82 & 38 & 6 \\
				Rooftop solar program & 76 & 51 & 34 \\
				Community solar project & 97 & 44 & 26 \\
				Postpone investment & 28 & 28 & 28 \\
				\bottomrule
			\end{tabular}
		\end{center}
		\vspace{-5mm}
		
		\textbf{Questions/Tasks.}
		\begin{enumerate}
			\CriterionTask
			{subsec:solar-energy-investment-c4}{4}
			{What would the optimistic decision be (Maximax)? What would the conservative decision be (Maximin)? \score{2}}
		\end{enumerate}
		
		\phantomsection
		\label{subsec:solar-energy-investment-c4}
		\subsection*{C4 - Maximax and Maximin \score{2}}
		
		\secInicio[sec:solutions-extraction][sec:solar-energy-investment-c4]
		
		\begin{multicols}{2}
			
			\textbf{Complete the table using letters first.}
			
			\[
			\begin{array}{c|ccccc}
				& A_1 & A_2 & A_3 & A_4 & \text{Decision}\\
				\hline
				\text{Max} &
				M_1 & M_2 & M_3 & M_4 &
				A^*_{\mathrm{Maximax}}=A_3
				\\
				\text{Min} &
				m_1 & m_2 & m_3 & m_4 &
				A^*_{\mathrm{Maximin}}=A_2
			\end{array}
			\]
			
			\textbf{Maximax}
			
			\[
			\begin{array}{ll}
				M_1 &= \max\{82,38,6\}=82\\
				M_2 &= \max\{76,51,34\}=76\\
				M_3 &= \max\{97,44,26\}=97\\
				M_4 &= \max\{28,28,28\}=28
			\end{array}
			\]
			
			\textbf{Maximin}
			
			\[
			\begin{array}{ll}
				m_1 &= \min\{82,38,6\}=6\\
				m_2 &= \min\{76,51,34\}=34\\
				m_3 &= \min\{97,44,26\}=26\\
				m_4 &= \min\{28,28,28\}=28
			\end{array}
			\]
			
			\columnbreak
			
			\textbf{Replace the letters with their values.}
			
			\[
			\begin{array}{c|ccccc}
				& A_1 & A_2 & A_3 & A_4 & \text{Decision}\\
				\hline
				\text{Max} &
				82 & 76 & 97 & 28 &
				A^*_{\mathrm{Maximax}}=A_3
				\\
				\text{Min} &
				6 & 34 & 26 & 28 &
				A^*_{\mathrm{Maximin}}=A_2
			\end{array}
			\]
			
			Therefore, the optimistic decision is
			\[
			A^*_{\mathrm{Maximax}}=A_3
			quad\text{(Community solar project)}
			\]
			
			and the conservative decision is
			\[
			A^*_{\mathrm{Maximin}}=A_2
			\quad\text{(Rooftop solar program)}.
			\]
			
		\end{multicols}
		
	}
	
	\sectionScored
	{Restaurant Expansion Strategy}
	{C4-\sml{2}}
	{sec:restaurant-expansion-c4}
	{
		A restaurant company must choose a strategy for expanding a new location.
		Customer demand could be high, medium, or low. All payoffs are in
		thousands of dollars.
			\begin{center}
			\renewcommand{\arraystretch}{1.3}
			\small
			\begin{tabular}{lccc}
				\toprule
				& \multicolumn{3}{c}{Customer demand} \\
				\cmidrule(lr){2-4}
				Alternative & High & Medium & Low \\
				\midrule
				Large restaurant & 78 & 36 & 12 \\
				Medium restaurant & 68 & 49 & 24 \\
				Small restaurant & 91 & 47 & 29 \\
				No expansion & 61 & 42 & 37 \\
				\bottomrule
			\end{tabular}
		\end{center}
		\vspace{-5mm}
		
		\textbf{Questions/Tasks.}
		\begin{enumerate}
			\CriterionTask
			{subsec:restaurant-expansion-c4}{4}
			{What would the optimistic decision be (Maximax)? What would the conservative decision be (Maximin)? \score{2}}
		\end{enumerate}
		
		\phantomsection
		\label{subsec:restaurant-expansion-c4}
		\subsection*{C4 - Maximax and Maximin \score{2}}
		
		\secInicio[sec:solutions-extraction][sec:restaurant-expansion-c4]
		
		\begin{multicols}{2}
			
			\textbf{Complete the table using letters first.}
			
			\[
			\begin{array}{c|ccccc}
				& A_1 & A_2 & A_3 & A_4 & \text{Decision}\\
				\hline
				\text{Max} &
				M_1 & M_2 & M_3 & M_4 &
				A^*_{\mathrm{Maximax}}=A_3
				\\
				\text{Min} &
				m_1 & m_2 & m_3 & m_4 &
				A^*_{\mathrm{Maximin}}=A_4
			\end{array}
			\]
			
			\textbf{Maximax}
			
			\[
			\begin{array}{ll}
				M_1 &= \max\{78,36,12\}=78\\
				M_2 &= \max\{68,49,24\}=68\\
				M_3 &= \max\{91,47,29\}=91\\
				M_4 &= \max\{61,42,37\}=61
			\end{array}
			\]
			
			\textbf{Maximin}
			
			\[
			\begin{array}{ll}
				m_1 &= \min\{78,36,12\}=12\\
				m_2 &= \min\{68,49,24\}=24\\
				m_3 &= \min\{91,47,29\}=29\\
				m_4 &= \min\{61,42,37\}=37
			\end{array}
			\]
			
			\columnbreak
			
			\textbf{Replace the letters with their values.}
			
			\[
			\begin{array}{c|ccccc}
				& A_1 & A_2 & A_3 & A_4 & \text{Decision}\\
				\hline
				\text{Max} &
				78 & 68 & 91 & 61 &
				A^*_{\mathrm{Maximax}}=A_3
				\\
				\text{Min} &
				12 & 24 & 29 & 37 &
				A^*_{\mathrm{Maximin}}=A_4
			\end{array}
			\]
			
			Therefore, the optimistic decision is
			\[
			A^*_{\mathrm{Maximax}}=A_3
			\quad\text{(Small restaurant)}
			\]
			
			and the conservative decision is
			\[
			A^*_{\mathrm{Maximin}}=A_4
			\quad\text{(No expansion)}.
			\]
			
		\end{multicols}
	}
	
	\sectionScored
	{Renewable Energy Investment Strategy}
	{C4-\sml{2}}
	{sec:renewable-energy-c4}
	{
		A renewable energy company must choose a strategy for introducing a new
		solar power system. Energy market conditions could be highly favourable,
		favourable, stable, or unfavourable. All payoffs are in thousands of dollars.
			\begin{center}
			\renewcommand{\arraystretch}{1.3}
			\small
			\begin{tabular}{lcccc}
				\toprule
				& \multicolumn{4}{c}{Energy market conditions} \\
				\cmidrule(lr){2-5}
				Alternative & Highly favourable & Favourable & Stable & Unfavourable \\
				\midrule
				Full-scale solar & 132 & 61 & 10 & -38 \\
				Expanded solar & 91 & 67 & 32 & 6 \\
				Basic solar & 65 & 54 & 40 & 22 \\
				Solar leasing & 51 & 46 & 37 & 30 \\
				\bottomrule
			\end{tabular}
		\end{center}
		\vspace{-5mm}
		
		\textbf{Questions/Tasks.}
		\begin{enumerate}
			\CriterionTask
			{subsec:renewable-energy-c4}{4}
			{What would the optimistic decision be (Maximax)? What would the conservative decision be (Maximin)? \score{2}}
		\end{enumerate}
		
		\phantomsection
		\label{subsec:renewable-energy-c4}
		\subsection*{C4 - Maximax and Maximin \score{2}}
		
		\secInicio[sec:solutions-extraction][sec:renewable-energy-c4]
		
		\begin{multicols}{2}
			
			\textbf{Complete the table using letters first.}
			
			\[
			\begin{array}{c|ccccc}
				& A_1 & A_2 & A_3 & A_4 & \text{Decision}\\
				\hline
				\text{Max} &
				M_1 & M_2 & M_3 & M_4 &
				A^*_{\mathrm{Maximax}}=A_1
				\\
				\text{Min} &
				m_1 & m_2 & m_3 & m_4 &
				A^*_{\mathrm{Maximin}}=A_4
			\end{array}
			\]
			
			\textbf{Maximax}
			
			\[
			\begin{array}{ll}
				M_1 &= \max\{132,61,10,-38\}=132\\
				M_2 &= \max\{91,67,32,6\}=91\\
				M_3 &= \max\{65,54,40,22\}=65\\
				M_4 &= \max\{51,46,37,30\}=51
			\end{array}
			\]
			
			\textbf{Maximin}
			
			\[
			\begin{array}{ll}
				m_1 &= \min\{132,61,10,-38\}=-38\\
				m_2 &= \min\{91,67,32,6\}=6\\
				m_3 &= \min\{65,54,40,22\}=22\\
				m_4 &= \min\{51,46,37,30\}=30
			\end{array}
			\]
			
			\columnbreak
			
			\textbf{Replace the letters with their values.}
			
			\[
			\begin{array}{c|ccccc}
				& A_1 & A_2 & A_3 & A_4 & \text{Decision}\\
				\hline
				\text{Max} &
				132 & 91 & 65 & 51 &
				A^*_{\mathrm{Maximax}}=A_1
				\\
				\text{Min} &
				-38 & 6 & 22 & 30 &
				A^*_{\mathrm{Maximin}}=A_4
			\end{array}
			\]
			
			Therefore, the optimistic decision is
			\[
			A^*_{\mathrm{Maximax}}=A_1
			\quad\text{(Full-scale solar)}
			\]
			
			and the conservative decision is
			\[
			A^*_{\mathrm{Maximin}}=A_4
			\quad\text{(Solar leasing)}.
			\]
			
		\end{multicols}
	}
			
	\sectionScored
	{Summarized Solutions}
	{C1-C4}
	{sec:summary}
	{
	\subsection{Meal Delivery Service \hfill C1--1 \quad C2--2}
	\begin{center}
		\begin{minipage}[t]{0.6\textwidth}
			\begin{tabularx}{\textwidth}
				{@{}p{1.5cm}>{\raggedright\arraybackslash}X@{}}
				\toprule
				Element & Description \\
				\midrule
				Alt. &
				\textbf{A1:} Standard delivery service \quad
				\textbf{A2:} Premium delivery service \\
				
				States &
				\textbf{S1:} High demand \quad
				\textbf{S2:} Low demand \\
				
				Events &
				Actual customer demand after the delivery service decision \\
				
				Cons. &
				Profit in thousand US dollars for each alternative--state pair \\
				\bottomrule
			\end{tabularx}
		\end{minipage}
		\hfill
		\begin{minipage}[t]{0.35\textwidth}
			\begin{center}
				\begin{tabular}{lcc}
					\toprule
					Alt. &  \(\mathbf{S_1}\) & \(\mathbf{S_2}\)\\
					\midrule
					\( \mathbf{A_1}\) & 20 & 12\\
					\(\mathbf{A_2}\) & 34 & 7\\
					\bottomrule
				\end{tabular}
			\end{center}
		\end{minipage}
	\end{center}
	
	\subsection{Small Coffee Roastery Production \hfill C1--1 \quad C2--3}
	\begin{center}
		\begin{minipage}[t]{0.55\textwidth}
			\begin{tabularx}{\textwidth}{
					@{}p{1.5cm}>{\raggedright\arraybackslash}X@{}
				}
				\toprule
				Element & Description \\
				\midrule
				Alt. &
				\textbf{A1:} Standard roasting machine \quad
				\textbf{A2:} Industrial roasting machine \\
				
				States &
				\textbf{S1:} High demand \quad
				\textbf{S2:} Low demand \\
				
				Events &
				Actual customer demand after the roasting machine decision \\
				
				Cons. &
				Profit in thousand BRL for each alternative--state pair \\
				\bottomrule
			\end{tabularx}
		\end{minipage}
		\hfill
		\begin{minipage}[t]{0.4\textwidth}
			\begin{tabular}{lcc}
				\toprule
				Alt. & \(\mathbf{S_1}\) & \(\mathbf{S_2}\)\\
				\midrule
				\( \mathbf{A_1}\)
				&
				\begin{tabular}{c}
					$18(4)+7$\\
					$=79$
				\end{tabular}
				&
				\begin{tabular}{c}
					$12(4)-4$\\
					$=44$
				\end{tabular}
				\\
				
				\( \mathbf{A_2}\)
				&
				\begin{tabular}{c}
					$15(6)+7$\\
					$=97$
				\end{tabular}
				&
				\begin{tabular}{c}
					$10(6)-4$\\
					$=56$
				\end{tabular}
				\\
				\bottomrule
			\end{tabular}
		\end{minipage}
	\end{center}
	
	\subsection{Online Fitness Membership \hfill C1--2 \quad C2--5}
	\begin{center}
		\begin{minipage}[t]{0.55\textwidth}
			\begin{tabularx}{\textwidth}{
					@{}p{1.5cm}>{\raggedright\arraybackslash}X@{}
				}
				\toprule
				Element & Description \\
				\midrule
				Alt. &
				\textbf{A1:} Standard membership \quad
				\textbf{A2:} Flexible membership \\
				
				States &
				\textbf{S1:} High customer demand \quad
				\textbf{S2:} Low customer demand \\
				
				Events &
				Actual customer demand after the membership is selected \\
				
				Cons. &
				Profit in thousand BRL for each alternative--state pair \\
				\bottomrule
			\end{tabularx}
		\end{minipage}
		\hfill
		\begin{minipage}[t]{0.4\textwidth}
			\begin{tabular}{lcc}
				\toprule
				Alt. & \(\mathbf{S_1}\) & \(\mathbf{S_2}\)\\
				\midrule
				\( \mathbf{A_1}\)
				&
				\begin{tabular}{c}
					$15(5+3)+5$\\
					$=125$
				\end{tabular}
				&
				\begin{tabular}{c}
					$9(5+1)+5$\\
					$=59$
				\end{tabular}
				\\
				
				\( \mathbf{A_2}\)
				&
				\begin{tabular}{c}
					$12(7+3)+8$\\
					$=128$
				\end{tabular}
				&
				\begin{tabular}{c}
					$7(7+1)+8$\\
					$=64$
				\end{tabular}
				\\
				\bottomrule
			\end{tabular}
			
		\end{minipage}
	\end{center}
	
	\subsection{Neighborhood Bakery Expansion \hfill C1--1 \quad C3--2}
	\begin{center}
		\begin{minipage}[t]{0.55\textwidth}
			\begin{tabularx}{\textwidth}{
					@{}p{1.5cm}>{\raggedright\arraybackslash}X@{}
				}
				\toprule
				Element & Description \\
				\midrule
				Alt. &
				\textbf{A1:} Introduce breakfast menu \quad
				\textbf{A2:} Maintain current selection \\
				
				States &
				\textbf{S1:} Strong traffic \quad
				\textbf{S2:} Weak traffic \\
				
				Events &
				Actual customer traffic after the product selection decision \\
				
				Cons. &
				Profit in million BRL for each alternative--state pair \\
				\bottomrule
			\end{tabularx}
		\end{minipage}
		\hfill
		\begin{minipage}[t]{0.4\textwidth}
			\begin{tabular}{lcc}
				\toprule
				Alt. & \(\mathbf{S_1}\) & \(\mathbf{S_2}\)\\
				\midrule
				\( \mathbf{A_1}\)
				&
				\begin{tabular}{c}
					$68-23$\\
					$=45$
				\end{tabular}
				&
				\begin{tabular}{c}
					$38-21$\\
					$=17$
				\end{tabular}
				\\
				
				\( \mathbf{A_2}\)
				&
				\begin{tabular}{c}
					$59-24$\\
					$=35$
				\end{tabular}
				&
				\begin{tabular}{c}
					$49-22$\\
					$=27$
				\end{tabular}
				\\
				\bottomrule
			\end{tabular}
			
		\end{minipage}
	\end{center}
	
	\subsection{Bakery Product Production \hfill C1--1 \quad C3--3}
	\begin{center}
		\begin{minipage}[t]{0.55\textwidth}
			\begin{tabularx}{\textwidth}{
					@{}p{1.5cm}>{\raggedright\arraybackslash}X@{}
				}
				\toprule
				Element & Description \\
				\midrule
				Alt. &
				\textbf{A1:} Produce cakes \quad
				\textbf{A2:} Produce pastries \\
				
				States &
				\textbf{S1:} High customer demand \quad
				\textbf{S2:} Low customer demand \\
				
				Events &
				Customer demand revealed after the production decision \\
				
				Cons. &
				Profit in thousand EUR for each product--demand-state pair \\
				\bottomrule
			\end{tabularx}
		\end{minipage}
		\hfill
		\begin{minipage}[t]{0.4\textwidth}
			\begin{tabular}{lcc}
				\toprule
				Alt. & \(\mathbf{S_1}\) & \(\mathbf{S_2}\)\\
				\midrule
				\( \mathbf{A_1}\)
				&
				\begin{tabular}{c}
					$9(20)-9(8)$\\
					$=108$
				\end{tabular}
				&
				\begin{tabular}{c}
					$5(20)-5(11)$\\
					$=45$
				\end{tabular}
				\\
				
				\( \mathbf{A_2}\)
				&
				\begin{tabular}{c}
					$6(16)-6(8)$\\
					$=48$
				\end{tabular}
				&
				\begin{tabular}{c}
					$4(16)-4(11)$\\
					$=20$
				\end{tabular}
				\\
				\bottomrule
			\end{tabular}
			
		\end{minipage}
	\end{center}
	
	
	\subsection{Meal Delivery Service Packages \hfill C1--2 \quad C3--5}
	{
		\begin{center}
			\begin{minipage}[t]{0.45\textwidth}
				\begin{tabularx}{\textwidth}{
						@{}p{1.5cm}>{\raggedright\arraybackslash}X@{}
					}
					\toprule
					Element & Description \\
					\midrule
					Alt. &
					\textbf{A1:} Offer a standard meal package \quad
					\textbf{A2:} Offer a premium meal package \\
					
					States &
					\textbf{S1:} High weekly demand \qquad \qquad \quad
					\textbf{S2:} Low weekly demand \\
					
					Events &
					Weekly market conditions revealed after the meal-package decision \\
					
					Cons. &
					Profit in thousand USD for each package--market-state pair \\
					\bottomrule
				\end{tabularx}
			\end{minipage}
			\hfill
			\begin{minipage}[t]{0.53\textwidth}
				\begin{tabular}{lcc}
					\toprule
					Alt. & \(\mathbf{S_1}\) & \(\mathbf{S_2}\)\\
					\midrule
					\( \mathbf{A_1}\)
					&
					\begin{tabular}{c}
						$10(7+3)-[10(3)+14]$\\
						$=56$
					\end{tabular}
					&
					\begin{tabular}{c}
						$7(7+1)-[7(3)+18]$\\
						$=17$
					\end{tabular}
					\\
					
					\( \mathbf{A_2}\)
					&
					\begin{tabular}{c}
						$8(12+3)-[8(5)+14]$\\
						$=66$
					\end{tabular}
					&
					\begin{tabular}{c}
						$5(12+1)-[5(5)+18]$\\
						$=22$
					\end{tabular}
					\\
					\bottomrule
				\end{tabular}
				
				
			\end{minipage}
		\end{center}
	}
	
	\newpage
	
	\begin{multicols}{3}
	\subsection{Solar Energy Investment Strategy \hfill C4--2}
	\[
	\begin{array}{c|cc}
		& \text{Max} & \text{Min}\\
		\hline
		A_1 & 82 & 6\\
		A_2 & 76 & 34\\
		A_3 & 97 & 26\\
		A_4 & 28 & 28\\
		\hline
		A^*_{\mathrm{Maximax}} &
		A_3 & \\
		A^*_{\mathrm{Maximin}} & & A_2
	\end{array}
	\]
	
	\subsection{Restaurant Expansion Strategy \hfill C4--2}
	\[
	\begin{array}{c|cc}
		& \text{Max} & \text{Min}\\
		\hline
		A_1 & 78 & 12\\
		A_2 & 68 & 24\\
		A_3 & 91 & 29\\
		A_4 & 61 & 37\\
		\hline
		A^*_{\mathrm{Maximax}} &
		A_3 & \\
		A^*_{\mathrm{Maximin}} & & A_4
	\end{array}
	\]
	
	\subsection{Renewable Energy Investment Strategy \hfill C4--2}
	\[
	\begin{array}{c|cc}
		& \text{Max} & \text{Min}\\
		\hline
		A_1 & 132 & -38\\
		A_2 & 91 & 6\\
		A_3 & 65 & 22\\
		A_4 & 51 & 30\\
		\hline
		A^*_{\mathrm{Maximax}} &
		A_1 & \\
		A^*_{\mathrm{Maximin}} & & A_4
	\end{array}
	\]
	\end{multicols}
	
\end{document}
